\subsection{数据预处理与特征分析}

\subsubsection{输入数据与口径}

四个问题的输入来自附件 1--4，其用途与读取口径如表~\ref{tab:data}。需要说明的是：附件 1 中的负载与光伏预测列只服务于问题一的“单日代表”语义，附件 3 的光伏预报只服务于问题三与链 \emph{4-3}，附件 4 的电价只服务于问题四。各问题的输入边界在建模时即已固定，不做跨问题借用，以保证四个问题之间互不“串答案”。

\begin{table}[htbp]
\centering
\small
\caption{输入数据及其在各问题中的用途}
\label{tab:data}
\begin{tabular}{p{1.8cm}p{3.4cm}p{7.0cm}}
\toprule
数据 & 结构 & 用途与读取口径 \\
\midrule
附件 1 & 145 行 $\times$ 4 列，10 分钟粒度单日 & 问题一至三的电价（逐日重复 365 次）；问题一的负载与光伏预测。电价取值于 $[0.3713,1.3952]$ 元/kWh，144 点之和为 110.3324 元/kWh \\
附件 2 & 365 行 $\times$ 144 列，实际值 & 问题二至四的负载与光伏实际功率。负载取值于 $[1995.7176,7978.8849]$ kW，光伏取值于 $[0,10216.20]$ kW \\
附件 3 & 365 天 $\times$ 4 个发布时刻 $\times$ 24 小时 & 问题三与链 \emph{4-3} 的整点光伏预报，需降尺度到 10 分钟 \\
附件 4 & 366 行 $\times$ 145 列，实时电价 & 问题四唯一的价格来源。取值于 $[0.0076,1.7936]$ 元/kWh，全期均值 0.7661976 元/kWh \\
模板文件 & 附件 5 的 5 个结果文件 & 决定表格结构与填报口径；模板中的数值不作为任何结论依据 \\
\bottomrule
\end{tabular}
\end{table}

附件 4 有三个值得注意的统计事实。第一，电价存在极端低值：全期最小值为 0.0076 元/kWh，低于 0.10 元/kWh 的时段仅 28 个，因此以最小电价构造的费用下界极为宽松（不足实际费用的 1.2\%），量级自检必须改用构造性上界。第二，深谷时段的 5 倍紧急购电价仅为 0.038 元/kWh，讨论惩罚机制时不能用全期均价外推。第三，附件 4 的\textbf{日内形状与附件 1 的电价曲线高度一致}：同时段跨日均值曲线与附件 1 曲线的相关系数为 0.9999999956（四位小数即为 1.0000）且极值位置相同，差异体现为日间水平漂移（日均价在 0.5378 至 0.9290 元/kWh 之间）叠加逐点扰动。这既解释了附件 1 电价能被逐日重复使用的原因，也说明在波动电价情景下采用“水平 $\times$ 时段”的双因子预测器有数据依据。

数据覆盖 2025 年 1 月 1 日至 12 月 31 日共 365 天、无缺日无重复；按结果文件要求，交付期为 2025 年 2 月 1 日至 12 月 31 日共 334 天，1 月作为储能预热期参与滚动递推但不计入交付。全部问题均自 2025 年 1 月 1 日 0:00 的 6000 kWh 起算，日间通过储电量连续性衔接，不施加逐日周期约束。

\subsubsection{负载与光伏的结构特征}

全年 52\,560 个时段中，出现光伏功率大于负载功率（即存在盈余）的时段有 10\,130 个，占 19.3\%。这类时段是储能充电的天然窗口，也是弃光可能发生的时段。交付期净负荷电量为 17\,939\,189.494967 kWh，该数值只与附件 2 有关，因此在问题二、问题三、链 \emph{4-2} 与链 \emph{4-3} 中\textbf{逐位相同}——它构成一个有用的交叉校验量：若某个模型算出的净负荷与此不符，说明取数或时间对齐出错。

由图~\ref{fig:input} 可见，负载在午间与傍晚出现双峰，光伏在正午达到最大，于是净负荷在正午被光伏压得很低、在傍晚急剧抬升。这意味着低价充电窗口集中在凌晨与正午、高价放电窗口集中在傍晚与夜间，正是分时电价套利的典型场景；而充放电之间的能量搬运要付出 19\% 的往返损耗，因此只有当价差足以覆盖损耗时套利才发生。

\begin{figure}[htbp]
\centering
\includegraphics[width=0.80\textwidth]{fig01_inputs.png}
\caption{附件 1 代表性单日的电价、小区负载与光伏预测功率，以及净负荷电量}
\label{fig:input}
\end{figure}

\subsubsection{异常与边界检查}

对全部输入做了四项检查：其一，缺失与重复检查——附件 2--4 无缺日无重复，附件 3 的日期列按前向填充还原为 365 天 $\times$ 4 个发布时刻；其二，量纲检查——功率乘以 $\Delta t$ 得电量，费用等于电价与电量之积且不乘 $\Delta t$；其三，时序一致性检查——附件 1 与附件 2 的 144 点采用同一左端点口径，若误用右端点会使全部时段整体错位 10 分钟；其四，净负荷一致性检查——四问共用同一净负荷数值。上述检查在每次计算后都以只读方式重新执行，检查脚本与计算脚本相互独立。
